% GENERATED FILE. Edit canonical lessons, then run npm run generate:books.
% canonical-source-hash: fnv1a64:fdb92fe91488f0dc
% canonical-lessons: RU-C345-spravedlivyy, RU-C345-nespravedlivyy, RU-C345-spravedlivost, RU-C345-razumnyy, RU-C345-istina

\chapter{Fair and unfair}
\label{ch:ru-fair-and-unfair}
\hlchaptermodality{345}

\begin{chapteropening}
\textbf{By the end of this chapter:} \emph{I can say whether something is fair or unfair, speak about justice, call an idea reasonable and speak of the truth.}

Complete the last lesson of chapter 345: I can say whether something is fair or unfair, speak about justice, call an idea reasonable and speak of the truth.
\end{chapteropening}

\section[\texorpdfstring{\ru{справедливый} (spravedlívyy)}{справедливый (spravedlívyy)}]{\ru{справедливый} (spravedlívyy) — fair, even-handed}

\label{lesson:RU-C345-spravedlivyy}



\begin{quote}
Before the new one: say the Russian for to reason, then the Russian for obvious.
\end{quote}

\subsection*{You'll want to know: \ru{справедливый}}
\textbf{\ru{справедливый}} (\emph{spravedlívyy}) — \textquotedblleft{}fair, even-handed\textquotedblright{}.

\begin{quote}
\ru{Наш} \ru{учитель} \ru{строгий}, \ru{но} \ru{справедливый}.
\end{quote}

\begin{quote}
\emph{Nash uchítel strógiy, no spravedlívyy.}
\end{quote}

\begin{quote}
Our teacher is strict but fair.
\end{quote}

\emph{\ru{Справедливо}} — fairly, or it is fair.

\begin{cousinweb}
\textbf{\ru{справедливый}} is usually traced to the root of \emph{\ru{правда}} \textquotedblleft{}truth\textquotedblright{} and \emph{\ru{право}} \textquotedblleft{}right\textquotedblright{}, Proto-Slavic \emph{*pravŭ} \textquotedblleft{}right, correct\textquotedblright{} (EtymDB), through a Church Slavonic formation.
\end{cousinweb}

\begin{grammarlens}[title={What you've built}]
One more word for telling, explaining and describing at B1.
\end{grammarlens}

\subsection*{Your turn}
\begin{itemize}
\raggedright
  \item \emph{Say it:} \emph{spravedlívyy}
  \item \emph{Say it:} \emph{spravedlívyy}, once more
  \item \emph{Say it:} a sentence of your own with \emph{spravedlívyy}
  \item \emph{Recall:} say the Russian for to reason, then the Russian for obvious, then say \emph{spravedlívyy} again
\end{itemize}

\begin{tcolorbox}[breakable,colback=teal!4,colframe=teal!35!black,title={Before you move on}]
What does \ru{справедливый} mean? (\textquotedblleft{}Fair, even-handed\textquotedblright{}.) Say it once more.
\end{tcolorbox}
\section[\texorpdfstring{\ru{несправедливый} (nespravedlívyy)}{несправедливый (nespravedlívyy)}]{\ru{несправедливый} (nespravedlívyy) — unfair}

\label{lesson:RU-C345-nespravedlivyy}



\begin{quote}
Before the new one: say the Russian for obvious, then the Russian for fair.
\end{quote}

\subsection*{You'll want to know: \ru{несправедливый}}
\textbf{\ru{несправедливый}} (\emph{nespravedlívyy}) — \textquotedblleft{}unfair\textquotedblright{}.

\begin{quote}
\ru{Это} \ru{несправедливо}: \ru{он} \ru{работал} \ru{больше} \ru{всех}.
\end{quote}

\begin{quote}
\emph{Éto nespravedlívo: on rabótal bólshe vsekh.}
\end{quote}

\begin{quote}
That is unfair: he worked more than anyone.
\end{quote}

Short form \emph{\ru{несправедлив}}. \emph{\ru{Несправедливо}} — it is unfair.

\begin{cousinweb}
\textbf{\ru{несправедливый}} is \emph{\ru{не}-} \textquotedblleft{}not\textquotedblright{} in front of \emph{\ru{справедливый}} \textquotedblleft{}fair\textquotedblright{}, which is usually traced to the root of \emph{\ru{право}}, Proto-Slavic \emph{*pravŭ} \textquotedblleft{}right, correct\textquotedblright{} (EtymDB).
\end{cousinweb}

\begin{grammarlens}[title={What you've built}]
One more word for telling, explaining and describing at B1.
\end{grammarlens}

\subsection*{Your turn}
\begin{itemize}
\raggedright
  \item \emph{Say it:} \emph{nespravedlívyy}
  \item \emph{Say it:} \emph{nespravedlívyy}, once more
  \item \emph{Say it:} a sentence of your own with \emph{nespravedlívyy}
  \item \emph{Recall:} say the Russian for obvious, then the Russian for fair, then say \emph{nespravedlívyy} again
\end{itemize}

\begin{tcolorbox}[breakable,colback=teal!4,colframe=teal!35!black,title={Before you move on}]
What does \ru{несправедливый} mean? (\textquotedblleft{}Unfair\textquotedblright{}.) Say it once more.
\end{tcolorbox}
\section[\texorpdfstring{\ru{справедливость} (spravedlívost)}{справедливость (spravedlívost)}]{\ru{справедливость} (spravedlívost) — fairness, justice}

\label{lesson:RU-C345-spravedlivost}



\begin{quote}
Before the new one: say the Russian for fair, then the Russian for unfair.
\end{quote}

\subsection*{You'll want to know: \ru{справедливость}}
\textbf{\ru{справедливость}} (\emph{spravedlívost}) — \textquotedblleft{}fairness, justice\textquotedblright{}.

\begin{quote}
\ru{Люди} \ru{хотят} \ru{справедливости}.
\end{quote}

\begin{quote}
\emph{Lyúdi khotyát spravedlívosti.}
\end{quote}

\begin{quote}
People want justice.
\end{quote}

Feminine. After \emph{\ru{хотеть}} in this sense it takes the genitive.

\begin{cousinweb}
\textbf{\ru{справедливость}} is built on \emph{\ru{справедливый}} \textquotedblleft{}fair\textquotedblright{} with the abstract suffix \emph{-\ru{ость}}; the root is that of \emph{\ru{право}}, Proto-Slavic \emph{*pravŭ} (EtymDB).
\end{cousinweb}

\begin{grammarlens}[title={What you've built}]
One more word for telling, explaining and describing at B1.
\end{grammarlens}

\subsection*{Your turn}
\begin{itemize}
\raggedright
  \item \emph{Say it:} \emph{spravedlívost}
  \item \emph{Say it:} \emph{spravedlívost}, once more
  \item \emph{Say it:} a sentence of your own with \emph{spravedlívost}
  \item \emph{Recall:} say the Russian for fair, then the Russian for unfair, then say \emph{spravedlívost} again
\end{itemize}

\begin{tcolorbox}[breakable,colback=teal!4,colframe=teal!35!black,title={Before you move on}]
What does \ru{справедливость} mean? (\textquotedblleft{}Fairness, justice\textquotedblright{}.) Say it once more.
\end{tcolorbox}
\section[\texorpdfstring{\ru{разумный} (razúmnyy)}{разумный (razúmnyy)}]{\ru{разумный} (razúmnyy) — reasonable, sensible}

\label{lesson:RU-C345-razumnyy}



\begin{quote}
Before the new one: say the Russian for unfair, then the Russian for fairness.
\end{quote}

\subsection*{You'll want to know: \ru{разумный}}
\textbf{\ru{разумный}} (\emph{razúmnyy}) — \textquotedblleft{}reasonable, sensible\textquotedblright{}.

\begin{quote}
\ru{Это} \ru{разумное} \ru{решение}.
\end{quote}

\begin{quote}
\emph{Éto razúmnoye reshéniye.}
\end{quote}

\begin{quote}
That is a sensible decision.
\end{quote}

Short form \emph{\ru{разумен}}. \emph{\ru{В} \ru{разумных} \ru{пределах}} — within reason.

\begin{cousinweb}
\textbf{\ru{разумный}} is built on \emph{\ru{разум}} \textquotedblleft{}reason\textquotedblright{}, which Russian borrowed from Church Slavonic; it comes from Proto-Slavic \emph{*orzumŭ}, built on \emph{*umŭ} \textquotedblleft{}mind\textquotedblright{}, the word that gives \emph{\ru{ум}} (EtymDB). The \emph{\ru{ра}-} at the start is the Church Slavonic shape where Russian also has \emph{\ru{роз}-}.
\end{cousinweb}

\begin{grammarlens}[title={What you've built}]
One more word for telling, explaining and describing at B1.
\end{grammarlens}

\subsection*{Your turn}
\begin{itemize}
\raggedright
  \item \emph{Say it:} \emph{razúmnyy}
  \item \emph{Say it:} \emph{razúmnyy}, once more
  \item \emph{Say it:} a sentence of your own with \emph{razúmnyy}
  \item \emph{Recall:} say the Russian for unfair, then the Russian for fairness, then say \emph{razúmnyy} again
\end{itemize}

\begin{tcolorbox}[breakable,colback=teal!4,colframe=teal!35!black,title={Before you move on}]
What does \ru{разумный} mean? (\textquotedblleft{}Reasonable, sensible\textquotedblright{}.) Say it once more.
\end{tcolorbox}
\section[\texorpdfstring{\ru{истина} (ístina)}{истина (ístina)}]{\ru{истина} (ístina) — the truth (in a deep sense)}

\label{lesson:RU-C345-istina}



\begin{quote}
Before the new one: say the Russian for fairness, then the Russian for reasonable.
\end{quote}

\subsection*{You'll want to know: \ru{истина}}
\textbf{\ru{истина}} (\emph{ístina}) — \textquotedblleft{}the truth (in a deep sense)\textquotedblright{}.

\begin{quote}
\ru{Он} \ru{думает}, \ru{что} \ru{знает} \ru{истину}.
\end{quote}

\begin{quote}
\emph{On dúmayet, chto znáyet ístinu.}
\end{quote}

\begin{quote}
He thinks he knows the truth.
\end{quote}

Feminine. More solemn than \emph{\ru{правда}}: a truth in philosophy or science.

\begin{cousinweb}
\textbf{\ru{истина}} is built on the old adjective \emph{\ru{истый}} \textquotedblleft{}true, genuine\textquotedblright{}; EtymDB sets Latvian \emph{īstens} \textquotedblleft{}real, true\textquotedblright{} beside it as a cousin (EtymDB).
\end{cousinweb}

\begin{grammarlens}[title={What you've built}]
One more word for telling, explaining and describing at B1.
\end{grammarlens}

\subsection*{Your turn}
\begin{itemize}
\raggedright
  \item \emph{Say it:} \emph{ístina}
  \item \emph{Say it:} \emph{ístina}, once more
  \item \emph{Say it:} a sentence of your own with \emph{ístina}
  \item \emph{Recall:} say the Russian for fairness, then the Russian for reasonable, then say \emph{ístina} again
\end{itemize}

\begin{tcolorbox}[breakable,colback=teal!4,colframe=teal!35!black,title={Before you move on}]
What does \ru{истина} mean? (\textquotedblleft{}The truth (in a deep sense)\textquotedblright{}.) Say it once more.
\end{tcolorbox}
